\documentclass[10pt]{ctexart}
\usepackage[a4paper,margin=1in]{geometry}
\usepackage{amsmath,amssymb,booktabs,array}
\usepackage{enumitem}
\usepackage{listings}
\usepackage{xcolor}
\usepackage{microtype}
\usepackage{hyperref}
\setlist[itemize]{leftmargin=1.5em}
\setlist[enumerate]{leftmargin=1.8em}
\lstset{basicstyle=\ttfamily\small,breaklines=true,columns=fullflexible,keepspaces=true,frame=single,framerule=0.2pt,xleftmargin=0.5em,xrightmargin=0.5em}
\title{014: StepPPO-v3 Method: Value Head on GiGPO}
\author{}
\date{2026-06-09}
\begin{document}
\maketitle
\vspace{-1.2em}

\section*{一句话定义}
当前分支里的 StepPPO-v3 就是一个 value-head-only 实验：在 GiGPO 训练上加一个 actor-side shared value head，训练它去拟合 step-level value/return，同时确认它不会影响 GiGPO 原本的 policy 训练。

这不是新的 policy algorithm，不默认引入 pure step advantage、action-level PPO ratio、\path{step_ppo_v3} estimator、separate critic 或额外 reward shaping。v3 的问题只有两个：
\begin{enumerate}
  \item value head 能不能稳定学出来；
  \item 加了 value head 后，GiGPO 的训练行为会不会变差。
\end{enumerate}

\section*{和 v2 的关系}
StepPPO-v2 中最值得单独验证的部分是 shared actor value head。v3 保留这个核心点，但不把 policy advantage 同时换掉，也不把 value loss 默认回传 actor backbone。这样做是为了隔离 value head 本身：如果 value head 都学不稳，或者会拖垮 GiGPO，那么后续更复杂的 StepPPO 设计没有意义。

\section*{训练对象}
给定一个 WebShop step sample，actor backbone 产生 hidden states，value head 在 action 前的 context boundary 上预测一个 scalar：
\begin{lstlisting}
b_i = full_sequence_length_i - response_length_i - 1
V_i = value_head(hidden_state_i[b_i])
\end{lstlisting}

主线默认：
\begin{lstlisting}
actor_rollout_ref.actor.value_head.enable=True
actor_rollout_ref.actor.value_head.detach_value_backbone=True
actor_rollout_ref.actor.value_head.value_loss_coef=0.03
actor_rollout_ref.actor.value_head.cliprange_value=0.5
\end{lstlisting}

\path{detach_value_backbone=true} 的含义是：value loss 只训练 value head，不通过 value regression 更新 actor backbone。GiGPO policy gradient 仍然按原路径更新 actor。

\section*{Policy 保持 GiGPO 不变}
v3 主线必须保持：
\begin{lstlisting}
algorithm.adv_estimator=gigpo
algorithm.gigpo.step_advantage_w=1.0
algorithm.gigpo.mode=mean_norm
actor_rollout_ref.actor.policy_loss.loss_mode=vanilla
\end{lstlisting}

也就是说，policy advantage、policy loss、rollout group 语义和 GiGPO baseline 保持一致。value head 是 auxiliary head，不参与当前 policy advantage 计算。

\section*{Value Target}
value head 的 target 使用 step-level GAE return。当前可复用 GiGPO-v1 value-aux trainer 中的路径：从 raw per-step reward 构造 value 用的 immediate reward，再调用 StepPPO-style GAE：
\begin{lstlisting}
R_i = GAE_return(reward_i, V_i, traj_uid, step_idx)
\end{lstlisting}

主线配置：
\begin{lstlisting}
algorithm.gigpo_v1_value_aux.value_target=gae_returns
algorithm.gigpo_v1_value_aux.step_gamma=0.95
algorithm.gigpo_v1_value_aux.step_lam=0.90
\end{lstlisting}

invalid action penalty 只能扣一次。建议保留 \path{value_step_rewards} 字段来标记 value target 使用的 reward 口径，避免和 GiGPO 自身 step reward 语义混淆。

\section*{Value Loss}
value loss 是 step scalar loss：
\begin{lstlisting}
values_clipped = old_values + clamp(current_values - old_values,
                                    -cliprange_value, cliprange_value)
value_loss = 0.5 * mean(max(
    (current_values - step_returns) ** 2,
    (values_clipped - step_returns) ** 2,
))
\end{lstlisting}

总 loss 仍然是 GiGPO policy loss 加一个 auxiliary value loss：
\begin{lstlisting}
loss = gigpo_policy_loss + value_loss_coef * value_loss
\end{lstlisting}

由于主线 detach backbone，auxiliary value loss 不应改变 actor backbone 的梯度方向。

\section*{必须记录的指标}
只看 \path{actor/value_loss} 不够，v3 必须记录：
\begin{lstlisting}
actor/value_loss
actor/value_clipfrac
actor/value_rmse
actor/value_mae
actor/value_explained_variance
actor/value_pred_mean
actor/value_pred_std
actor/value_return_mean
actor/value_return_std
actor/value_loss_coef
\end{lstlisting}

同时必须和 GiGPO baseline 对比 policy 行为：
\begin{lstlisting}
episode/valid_action_ratio
response_length/mean
response_length/clip_ratio
actor/pg_loss
actor/grad_norm
val/success_rate
val/webshop_task_score (not success_rate)
\end{lstlisting}

\section*{成功标准}
v3 成功不要求立刻超过 GiGPO。最低标准是：
\begin{enumerate}
  \item value RMSE/MAE 稳定，不随训练发散；
  \item value explained variance 从接近 0 向正值移动，说明 value head 学到有用 signal；
  \item 加 value head 后，GiGPO 的 valid action ratio、response length、clip ratio 和 validation score 不明显回退；
  \item detach=true 主线不复现 detach=false value-aux 那种 validation collapse；
  \item actor grad norm 不因为 value loss 明显异常。
\end{enumerate}

\section*{实验脚本}
主线脚本：
\begin{lstlisting}
exps/run_webshop_step_ppo_v3_value_head_4gpu_paper_align_simple.sh
\end{lstlisting}

这个脚本名保留 StepPPO-v3，是研究命名；实现含义就是 GiGPO + detached value head。

\section*{不属于当前 v3 的内容}
以下内容不要写成 v3 默认设计：
\begin{itemize}
  \item 新增 \path{algorithm.adv_estimator=step_ppo_v3}；
  \item pure step advantage 替代 GiGPO advantage；
  \item action-level PPO ratio 或 sequence geometric ratio；
  \item separate critic model；
  \item value loss 默认回传 actor backbone；
  \item 一边改 policy，一边评估 value head 是否可行。
\end{itemize}

\end{document}
